\section{How we know what we know}

We have often quoted Evola's claim that he thinks like sound and well-born men\footnote{\url{https://gornahoor.net/?p=1026}} prior to 1789. As an example, we have referred to Girolamo Savonarola\footnote{\url{https://gornahoor.net/?p=2520}}, and it will be helpful to bring in more examples. (Of course, not all men were sound thinkers nor well-born, but those who were held positions of intellectual leadership.) Tonight, we refer to a passage from On the Love of God (1616) by \textbf{St. Francis de Sales}, describing how we know.

A philosopher may object to his perfunctory treatment, but the Saint was writing for those who would have accepted it as a matter of course. While logical and philosophical discussion of metaphysical ideas has its place, it is imperative to keep in mind that metaphysics is not a branch of philosophy. Physics is the science of sensible things; metaphysics is the science of states of consciousness. There is a plethora of turgid writings in journals and on the Internet that attempt to describe Tradition from the outside, using modernist ideas such as comparative religion, sociological theories, archeology and so on. However, truth is in the interiority of man (\textbf{St. Augustine}), so any return to Tradition is dependent on a change of consciousness (what Guenon calls an intellectual conversion\footnote{\url{https://gornahoor.net/?p=262}}).

So in following this passage, we must not adopt it as a theory among other theories. Even less, is it something that we debate or argue about. No, the only recourse is to observe or try to recreate the steps in our own consciousness; only then can we judge its merits. St. Francis describes the process of sense perception, which even a modernist can accept.

\begin{quotex}
When we look at anything, even though it is present before us, it is not itself united to our eyes, but merely sends them a certain representation or picture of itself, which is called its sensible species, by means of which we see the object.
\end{quotex}

But it is the next step that modernity fails to understand. \textbf{Immanuel Kant}, for example, got lost at this step when he realized we cannot know the thing-in-itself through sense perception alone. The scientist will attempt to formulate laws using words and mathematical models that describe sensual experience, which he then calls ``knowledge''. Those are both dead ends: in the former, we are reduced to perpetual ignorance. In the latter, we accept theories provisionally since they are always subject to change. Furthermore, an honest appraisal will reveal how little of our lives, and none of the most important parts of them, are amenable to science. The Doctor of the Church continues:

\begin{quotex}
When we contemplate or understand a thing, that which we understand is not united to our intellect except by means of another representation and most delicate spiritual image, which is called the intelligible species.
\end{quotex}


This means that alongside the image created by sense impressions, there is another image in thought. Thus, our mind is also an organ of perception, wherein it experiences things of the spirit (ideas). When the two images coincide, there is knowledge of the thing. St. Francis describes this process of understanding:

\begin{quotex}
Still, by how many routes and changes do these species get to our intellect. They arrive at the exterior senses, thence pass to the interior sense, then to the imagination, then to the active intellect, and come at last to the passive intellect. Hence, by passing through so many strainers and under so many files they may be purified, subtilized and perfected, and from sensible objects they become intelligible.
\end{quotex}


The sense object creates a reflection in the mind (imagination). The passive understanding means that we have the potential to understand, but the active understanding, or the effort of the mind to understand, is also required to make the potential actual. Now, in ordinary life, this happens instantly so we miss out on the entire process, except, perhaps, in recollecting the experience. But sometimes, say in poor lighting, we can observe the mind struggling to make out an object, when suddenly there is the moment of insight when we grasp it.

This all makes more sense when we apply it to human and purely spiritual knowledge. We have discussed the latter on many occasions, for example, in the discussion on St. Augustine\footnote{\url{https://gornahoor.net/?p=2842}}. St. Paul uses the analogy of a mirror for spiritual perception, darkly because our minds need be ``purified, subtilized, and perfected.'' But St. Francis goes on to say:

\begin{quotex}
But in heaven ... divinity will unite itself to our intellect without mediation of any species or representation whatsoever.
\end{quotex}


But this is exactly how we have been defining gnosis: direct, unmediated knowledge of spiritual truths. For the Hermetist, his goal is gnosis in this life, for which he purifies his consciousness\footnote{\url{https://www.meditationsonthetarot.com/moral-purification-of-the-will}}.

In between sensible things and intelligible things, there is the human world, constructed by thought, strained and defiled, which is the ground for both the greater and lesser battles. How we come to an understanding of this human world will be the topic of the next post in this series.

All quotes taken from Book 3, Chapter 11.


\flrightit{Posted on 2011-09-30 by Cologero}

\begin{center}* * *\end{center}

\begin{footnotesize}
\begin{sffamily}
\texttt{EXIT on 2011-09-30 at 10:03 said:} 

If I have ``merely repeated what every perennialist has written down since Guenon'', then that is a good thing, since innovation is not my goal. This post was to point out how a particular man thought in 1616, centuries \emph{before} Guenon. To read more from Exit, please visit Utopian Race. --- Cologero The article doesn't answer the question in the title. You've merely repeated what every perennialist has written down since Guenon. That doesn't mean you've attained metaphysical states of consciousness. Furthermore, when you speak of self-initiation that only confirms that you haven't attained those states. Or if you've received any other initiations please list them now. Otherwise you can't say you truly know metaphysics, and therefore your knowledge of it is merely philosophical.

\hfill

\texttt{Cologero on 2011-10-06 at 20:51 said:}

We at Gornahoor choose to discuss what Guenon does not, although he is probably correct that it is unprofitable. We also disagree with G. about the lack of traces. A message to that handful of readers who obsess over my personal being: \emph{You know who's not happy with their lives, when their busy discussing yours}.

\begin{quote}\footnotesize

Westerners, including even those who were true metaphysicians up to a certain point, have never known metaphysic in its entirety. Nevertheless there may have been individual exceptions, for nothing in principle precludes the existence, at any period or in any country, of men capable of attaining to complete metaphysical knowledge; and this would even be possible in the present Western world ... it is only right to point out that if such exceptions ever did exist, no written evidence has been found to prove it, nor do the generally known facts reveal traces of their influence; this absence of direct evidence however proves nothing in a negative sens, not should it occasion any surprise, for if cases of this sort did in fact occur, it can only have been thanks to very special circumstances of a nature that we cannot profitably discuss here. 
\flright{\textsc{Rene Guenon}, \textit{Introduction to the Study of Hindu Doctrines}}
\end{quote}


\hfill

\texttt{EXIT on 2011-10-09 at 08:33 said:}

My only interest in your personal being is the fact that upon reading Guenon you have taken his personality and key phrases but have only managed to misrepresent him since you have never attained metaphysical states of conscious with self-initiation which he severely denounced with good reason (since one lacking the spiritual influence cannot give it to oneself).

\hfill

\texttt{Cologero on 2011-10-09 at 12:10 said:}

\begin{quote}\footnotesize
We are all naturally more ready to censure errors than to praise things well done; and many men, from a kind of \textbf{innate malice}, and \emph{even when they clearly see the good}, strive with all effort and care to discover some fault or at least something that seems a fault.
\end{quote}


From \textbf{Baldesar Castiglione}, a wise Westerner.

Moral of the story: In the Comments section, please address specific points rather than general expressions of error, disagreement, etc.

\hfill

\texttt{Perennial on 2011-10-09 at 22:41 said:}

I must confess Mr. EXIT that I continue to be befuddled by your absence of citations or quotes. I do not necessarily fault your interpretation of Guenon, but I do fault you for not sharing your source that people might see your argument in his words and judge for themselves. I personally am more familiar with Evola then with Guenon, so citations would be a big help.


\hfill

\end{sffamily}
\end{footnotesize}

